\documentclass{elsarticle}
\usepackage{placeins}
\usepackage{amsmath,amssymb}
\usepackage{color}
\usepackage{fullpage}
\usepackage{booktabs}
\usepackage[english]{babel}
\usepackage{graphicx}
\usepackage{xcolor}
\usepackage[unicode]{hyperref}
\usepackage{float}
\usepackage{multirow}
\usepackage{enumitem}
\usepackage{appendix}
\usepackage{algorithm}
\usepackage{algpseudocode}
\usepackage{caption}
\usepackage{subcaption}
\usepackage{xspace}
\usepackage{upgreek}
\usepackage{mathtools}
\usepackage{setspace}
\doublespacing

\newcommand{\eg}{\emph{e.g.}\xspace}
\newcommand{\ie}{\emph{i.e.}\xspace}
\newcommand{\online}{\emph{online}\xspace}
\newcommand{\offline}{\emph{offline}\xspace}
\newcommand{\nick}[1]{\textcolor{blue}{#1}}
\newcommand{\rai}[1]{{\color{red}#1}}


\renewcommand{\vec}[1]{\boldsymbol{#1}}
\newcommand{\ten}[1]{\boldsymbol{#1}}
\newcommand{\mat}[1]{\mathsf{#1}}
\newcommand{\params}{\boldsymbol{\mu}}
\newcommand{\paramspace}{\boldsymbol{\Theta}}
\newcommand{\newparams}{\boldsymbol{\mu}^*}

\newcommand{\wi}{\mathrm{Wi}}

\newcommand{\meanof}[1]{\langle #1 \rangle}
\newcommand{\stdof}[1]{\mathrm{std}(#1)}

\newcommand{\subs}{\mathrm{s}}
\newcommand{\subp}{\mathrm{p}}
\newcommand{\subexp}{\mathrm{exp}}
\newcommand{\subbias}{\mathrm{bias}}
\newcommand{\submesh}{\mathrm{mesh}}        % NEW: mesh discretization error

\newcommand{\obs}{\mathrm{obs}}
\newcommand{\rom}{\mathrm{rom}}
\newcommand{\fom}{\mathrm{fom}}

\newcommand{\figWfull}{\linewidth}
\newcommand{\figWhalf}{0.49\linewidth}
\newcommand{\figWthird}{0.33\linewidth}

\begin{document}
\begin{frontmatter}

\title{Bayesian inference using extrudate swell of viscoelastic fluid}

\author[1]{Parajal Rai\corref{cor1}}
\ead{p.rai@tue.nl}
\author[1]{Michelle M.A. Spanjaards}
\author[1]{Ye Wang}
\author[1]{Patrick D. Anderson}
\author[1]{Nick Jaensson}
\cortext[cor1]{Corresponding author.}
\address[1]{Department of Mechanical Engineering,
            Eindhoven University of Technology,
            5600 MB Eindhoven, The Netherlands}


\begin{abstract}

\end{abstract}

\begin{keyword}
Bayesian inference \sep
extrudate swell \sep
reduced model \sep
viscoelasticity \sep
Markov chain Monte Carlo
\end{keyword}

\end{frontmatter}


\begin{abstract}
Extrudate swell is a sensitive manifestation of viscoelastic stress relaxation
and therefore provides a natural observable for material parameter inference.
However, direct Bayesian inference with a high-fidelity extrudate-swell
full-order model (FOM) is computationally prohibitive because each likelihood
evaluation requires solving a nonlinear free-surface viscoelastic flow problem.
This work develops a Bayesian inverse framework in which a non-intrusive
reduced-order model (ROM) is used as the forward model for the free-surface
profile. The ROM is trained from FOM snapshots of the fully swollen extrudate
surface using a proper orthogonal decomposition basis and either radial basis
function or Gaussian-process regression for the reduced coefficients. The
resulting online model maps constitutive parameters directly to the observed
curve, denoted here as \texttt{curve4}, and is inexpensive enough for Markov
chain Monte Carlo sampling. The framework is formulated for three common
viscoelastic constitutive models: Oldroyd-B, for which the inferred parameters
are the relaxation time and solvent-viscosity ratio; Giesekus, for which the
mobility parameter is additionally inferred; and Phan-Thien--Tanner (PTT), for
which the extensibility parameter replaces the Giesekus mobility parameter. In
addition to FOM-generated synthetic observations, we describe a video-based
pipeline in which a Blender rendering mimics experimental images, the fully
swollen free-surface profile is extracted, and the resulting curve is used for
Bayesian inference. The method connects high-fidelity ALE finite-element
simulation, non-intrusive ROM construction, computer-vision extraction of
extrudate profiles, and uncertainty-aware parameter estimation.
\end{abstract}

\noindent\textbf{Keywords:} extrudate swell; Bayesian inference; viscoelastic
fluid; Oldroyd-B; Giesekus; PTT; reduced-order model; full-order model;
free-surface flow; synthetic experiment

\section{Introduction}

\section{Bayesian inference framework}
\label{sec:bayes}

A key goal of this study is to infer the rheological parameters of a viscoelastic fluid from the geometry of the extrudate produced at the die exit. The free-surface profile serves as a natural and experimentally accessible observable, providing valuable information for constitutive parameter estimation.

Bayesian inference provides a formulation in which the unknown parameters are
modeled as random variables described by probability distributions, and the
discrepancy between the model and the data can be explicitly decomposed into
experimental noise, model bias, and numerical error. In this work, the Bayesian
framework is used exclusively for parameter identification, \ie we do not use it
to perform model selection. The forward model is the non-intrusive reduced-order
model constructed in Section~\ref{sec:rom}, which maps a constitutive parameter
vector directly onto the fully swollen free-surface profile, and the inference
returns the posterior density of the constitutive parameters given the observed
profile. The error model, prior distributions, likelihood, and Markov chain
Monte Carlo (MCMC) sampler are specified in Section~\ref{sec:bayesian_inference}.

% \subsection{Observation model}
% \label{subsec:obs_model}

We consider an optical measurement of the extrudate, in which the free surface is recorded in the plane containing the flow axis and the
resulting edge is converted into a height profile. The observation is taken in
the fully swollen state, \ie once the free surface has reached steady state and
no longer evolves in the streamwise direction, so that the data carry no
explicit time dependence. The profile is sampled at $N_h$ fixed axial stations
$x_{1},\dots,x_{N_h}$ downstream of the die exit, and the measurement is
written as the observation vector
\begin{equation}
\vec{h}_{\obs}
=
[\,h_{\obs,1},\dots,h_{\obs,N_h}\,] \in \mathbb{R}^{N_h},
\label{eq:observation_vector}
\end{equation}
where $h_{\obs,\,i} \in \mathbb{R}$ is the measured height of the free surface
at the axial position $x_{i}$, non-dimensionalized by the die half-height $H$.

Let ${d}_{i}(\params) \in \mathbb{R}$ denote the prediction of the forward
model at $x_{i}$ for the parameter vector $\params \in \paramspace$, where
$\paramspace$ is the parameter space. For an Oldroyd-B fluid
$\params=(\lambda,\beta)^{\top}$, whereas the Giesekus and PTT models introduce
one additional nonlinear parameter, as detailed in
Section~\ref{sec:bayesian_inference}. The forward model is the deterministic
mapping $\vec{d}: \paramspace \to \mathbb{R}^{N_h}$ that returns the predicted
free-surface heights at the $N_h$ sampling stations,
\begin{equation}
\vec{d}(\params)
=
[\,{d}_{1}(\params),\dots,{d}_{N_h}(\params)\,] \in \mathbb{R}^{N_h}.
\label{eq:forward_model}
\end{equation}
Evaluating $\vec{d}$ with the full-order model of Section~\ref{sec:fom} would
require the solution of a nonlinear free-surface viscoelastic flow problem for
every candidate parameter, which is prohibitive inside an MCMC loop. The
reduced-order model of Section~\ref{sec:rom} therefore takes its place, at the
cost of introducing a surrogate error, which is quantified in
Section~\ref{subsec:error_model}.

% ---- REVISED: observation model extended with ROM and mesh error terms -------
Kennedy and O'Hagan~\cite{kennedy2001calibration} decompose the difference
between an observation and a computer model into measurement noise and model
discrepancy. We extend this decomposition by two numerical error terms and
write the observations as
\begin{equation}
\vec{h}_{\obs}
\;=\;
\vec{d}(\params)
+
\vec{\varepsilon}_{\rom}(\params)
+
\vec{\varepsilon}_{\submesh}
+
\vec{\varepsilon}_{\subbias}
+
\vec{\varepsilon}_{\subexp},
\label{eq:error_decomposition}
\end{equation}
where $\vec{\varepsilon}_{\subexp} \in \mathbb{R}^{N_h}$ denotes the
experimental noise, $\vec{\varepsilon}_{\subbias} \in \mathbb{R}^{N_h}$ the
model bias (also called model error or model discrepancy),
$\vec{\varepsilon}_{\rom}(\params) \in \mathbb{R}^{N_h}$ the error of the
reduced-order model with respect to the full-order model it emulates, and
$\vec{\varepsilon}_{\submesh} \in \mathbb{R}^{N_h}$ the discretization error of
the full-order model. The four terms have distinct origins and distinct
structure. The experimental noise arises from the imaging chain, \eg sensor
noise and the uncertainty in locating the edge of the extrudate, and is
reasonably modeled as uncorrelated between sampling stations. The model bias
represents the inadequacy of the constitutive model used in the forward model,
\eg an Oldroyd-B model calibrated against the profile of a shear-thinning
fluid, together with any systematic distortion introduced when the profile is
extracted from an image. Because this error varies smoothly along the
extrudate, the bias is correlated between neighbouring stations, and it is this
correlation in $x$ that distinguishes it from the measurement noise. The ROM
error collects the truncation of the reduced basis and the regression error of
the surrogate for the reduced coefficients, and therefore depends on $\params$.
The mesh error accounts for the finite resolution of the finite-element
solutions from which the reduced-order model is built. A real measurement
carries no discretization error, so that all of it resides in the forward model.

The two numerical error terms are not inferred from the observed profile,
because they would be confounded with the model bias. Since the full-order
model is available, they are instead quantified beforehand. The mesh error is
obtained from a mesh convergence study, and its covariance is fixed before the
sampling. The ROM error is obtained from leave-one-out predictions of the
training snapshots; it is found to be much smaller than the other three
contributions and is neglected in the inference
(Section~\ref{subsec:error_model}). The amplitudes of the experimental noise and
of the model bias are inferred jointly with $\params$; the bias term is omitted
when the data are generated with the same constitutive model as the forward
model.
% ---- end of revised block ----------------------------------------------------

Given the observation model~\eqref{eq:error_decomposition}, the parameter
identification problem is posed in the Bayesian setting. Bayes' theorem
expresses the probability density of $\params$ given the observed free-surface
profile $\vec{h}_{\obs}$, referred to as the posterior, as
\begin{equation}
p(\params \mid \vec{h}_{\obs})
\;=\;
\frac{p(\vec{h}_{\obs} \mid \params)\,p(\params)}
     {p(\vec{h}_{\obs})}
\;=\;
\frac{\mathcal{L}(\params \mid \vec{h}_{\obs})\,p(\params)}
     {p(\vec{h}_{\obs})},
\label{eq:bayes_intro}
\end{equation}
where
$\mathcal{L}(\params \mid \vec{h}_{\obs}) \equiv p(\vec{h}_{\obs} \mid \params)$
is the likelihood, $p(\params)$ the prior, and
\begin{equation}
p(\vec{h}_{\obs})
\;=\;
\int_{\paramspace}
\mathcal{L}(\params \mid \vec{h}_{\obs})\,p(\params)\,
\mathrm{d}\params,
\label{eq:evidence}
\end{equation}
the evidence.

The prior $p(\params)$ encodes information available before the profile is
observed, such as positivity of the relaxation time, the physical restriction
of the viscosity ratio to the unit interval, expert knowledge, independent
rheometric measurements, or the range of parameters over which the reduced-order
model has been trained and is therefore valid. The likelihood
$\mathcal{L}(\params \mid \vec{h}_{\obs})$ quantifies how plausibly a parameter
$\params$ explains the observed profile. Through the observation
model~\eqref{eq:error_decomposition}, the measurement noise, the model bias,
and the mesh discretization error all enter the inference via this term.

The evidence $p(\vec{h}_{\obs})$ normalizes the posterior over $\paramspace$,
but its evaluation requires an integral over the entire parameter domain, which
becomes computationally prohibitive in high dimensions. The posterior is
therefore characterized using Markov chain Monte Carlo sampling, which generates
samples from $p(\params \mid \vec{h}_{\obs})$ using only ratios of posterior
densities, in which the evidence cancels.

The posterior distribution combines prior knowledge with information provided by the data through the likelihood. When the data are sparse or noisy, an informative prior can strongly influence the posterior, whereas with highly informative data, the likelihood becomes dominant. From the posterior, we obtain point estimates, credible intervals, and parameter correlations.

\section{Full Order Model}
\label{sec:fom}

The full-order model (FOM) is used to generate the high-fidelity snapshot data for the reduced-order model. It resolves the coupled viscoelastic flow, free-surface deformation, and mesh motion in the extrudate-swell problem.
This section summarizes the computational domain, governing equations, constitutive model, ALE formulation, and numerical discretization used in the FOM.

\subsection{Problem description}

\begin{figure}[!ht]
\centering
\centering
\includegraphics[width=0.5\textwidth]{Picture14.png}
\caption{2D schematic of extrudate-swell configuration.}
\label{geometry}

\end{figure}

We consider the extrusion of a viscoelastic fluid through a 2D axisymmetric into an open domain. As the fluid exits the die, the elastic stresses accumulated during the flow inside the die then relax, leading to an expansion of the free surface, commonly referred to as the \emph{extrudate-swell} phenomenon. A 2D schematic representation of the computational domain is shown in Fig.~\ref{geometry}. The die has a height $H = 1$ and a length $L_{\mathrm{die}} = 3$, followed by an open extrusion region of length $L_{\mathrm{extr}} = 5$ where the free surface evolves. The flow is driven by a prescribed flow rate $Q$, and the rheological behavior of the fluid is described using the Giesekus constitutive model.


\subsection{Balance and constitutive equations}

We consider an incompressible flow in the creeping-flow regime, where
inertial effects are negligible and no external body forces are present.
Under these assumptions, the governing equations reduce to the conservation
of linear momentum and mass:
\begin{align}
\nabla \cdot \ten{\sigma} &= \vec{0},
\label{eq:momentum_balance}
\\
\nabla \cdot \vec{u} &= 0.
\label{eq:mass_balance}
\end{align}
Here, $\vec{u}$ denotes the fluid velocity and $\ten{\sigma}$ is the
Cauchy stress tensor. The total stress is decomposed into an isotropic
pressure contribution, a Newtonian solvent stress, and a polymeric extra
stress:
\begin{equation}
\ten{\sigma}
=
-p\ten{I}
+
2\eta_{\mathrm{s}}\ten{D}
+
\ten{\tau},
\label{eq:stress_decomposition}
\end{equation}
where $p$ is the pressure, $\eta_{\mathrm{s}}$ is the solvent viscosity,
and $\ten{\tau}$ is the polymeric extra-stress tensor. The
rate-of-deformation tensor is defined as
\begin{equation}
\ten{D}
=
\frac{1}{2}
\left(
\nabla \vec{u}
+
\nabla \vec{u}^{T}
\right).
\end{equation}

The polymeric stress is related to the conformation tensor $\ten{c}$ through
\begin{equation}
\ten{\tau}
=
G\left(\ten{c}-\ten{I}\right),
\label{eq:poly_stress}
\end{equation}
where $G$ is the polymer shear modulus. The polymer viscosity is consequently
given by
\begin{equation}
\eta_{\mathrm{p}} = G\lambda,
\label{eq:polymer_viscosity}
\end{equation}
where $\lambda$ is the characteristic polymer relaxation time.

The evolution of the conformation tensor is written in the general form
\begin{equation}
\frac{D\ten{c}}{Dt}
=
\ten{L}\cdot\ten{c}
+
\ten{c}\cdot\ten{L}^{T}
-
f(\ten{c}),
\label{eq:generic_model}
\end{equation}
where $\ten{L}=\nabla\vec{u}$ is the velocity-gradient tensor. Equivalently,
using the upper-convected derivative,
\begin{equation}
\overset{\triangledown}{\ten{c}}
=
\frac{D\ten{c}}{Dt}
-
\ten{L}\cdot\ten{c}
-
\ten{c}\cdot\ten{L}^{T}.
\end{equation}

For the Giesekus model, the relaxation function is given by
\begin{equation}
f_{\mathrm{G}}(\ten{c})
=
\frac{1}{\lambda}
\left[
\left(\ten{c}-\ten{I}\right)
+
\alpha
\left(\ten{c}-\ten{I}\right)^2
\right],
\label{eq:giesekus_model}
\end{equation}
where $\alpha$ is the Giesekus mobility parameter. The quadratic term
introduces anisotropic drag and produces nonlinear rheological behavior,
including shear thinning.

For the exponential Phan--Thien--Tanner model, the relaxation function is
defined as
\begin{equation}
f_{\mathrm{PTT}}(\ten{c})
=
\frac{1}{\lambda}
\exp\left[
\epsilon\,
\operatorname{tr}
\left(
\ten{c}-\ten{I}
\right)
\right]
\left(
\ten{c}-\ten{I}
\right),
\label{eq:exponential_ptt_model}
\end{equation}
where $\epsilon$ is the PTT nonlinearity parameter.

The Giesekus and exponential PTT models reduce to the Oldroyd-B
model in the limits $\alpha=0$ and $\epsilon=0$, respectively.


\subsection{Boundary conditions}

We exploit the symmetry conditions about the channel centerline, and only the upper half of the computational domain is simulated. A fully developed velocity profile is imposed at the inlet $\Gamma_{\mathrm{in}}$, while no-slip conditions are applied along the die wall $\Gamma_{\mathrm{wall}}$. At the downstream boundary $\Gamma_{\mathrm{out}}$, a traction-free condition, together with $u_y=0$, is imposed. Similarly, symmetry boundary conditions are applied to $\Gamma_{\mathrm{sym}}$.

The free surface $\Gamma_{\mathrm{free}}$ is assumed to be traction-free, and the evolution of the free surface $\Gamma_{\mathrm{free}}$ is governed by the height function $h(x,t)$:
\begin{equation}
\frac{\partial h}{\partial t}
+
u_x \frac{\partial h}{\partial x}
=
u_y,
\qquad
x \in [0,L_{\mathrm{extr}}],
\label{eq:height_2d}
\end{equation}
The height function is initialized as $h(x,0)=H$ and  $h(0,t)=H$.


\subsection{Numerical methods}

The governing equations are discretized using the finite element method. Quadratic ($P_2$) elements are employed for the velocity field, while linear ($P_1$) elements are used for both the pressure and conformation tensor fields, resulting in a $P_2$–$P_1$–$P_1$ discretization. To account for the evolution of the free surface, an Arbitrary Lagrangian–Eulerian (ALE) formulation is employed~\cite{hughes1981ale,donea1982ale}.

To enhance numerical stability at moderate and high Weissenberg numbers, several stabilization techniques are employed. These include the log-conformation formulation~\cite{hulsen2005flow,fattal2004constitutive}, streamline-upwind/Petrov–Galerkin (SUPG) stabilization~\cite{brooks1982streamline}, and the DEVSS-G method~\cite{bogaerds2004stability}. Time integration is performed using a second-order backward differentiation scheme (BDF2).


\section{Reduced Order Model}
\label{sec:rom}

The extrudate-swell free surface evolves differently for different
rheological parameters. Let $\mat{h}(\params) \in \mathbb{R}^{N_h}$ denote the vector of
free surface curves for a parameter vector
$\params \in \mathcal{D} \subset \mathbb{R}^{d}$, where $N_h$ is the number of sampling points along the free surface. For a finite training
set
\[
    \Theta = \{\params_1,\dots,\params_M\} \subset \mathcal{D},
\]
the snapshots are assembled column-wise into
\begin{equation}
    \mat{X}
    =
    \bigl[
        \mat{h}(\params_1),
        \mat{h}(\params_2),
        \dots,
        \mat{h}(\params_M)
    \bigr]
    \in \mathbb{R}^{N_h \times M}.
    \label{eq:snapshot_matrix}
\end{equation}

The column space of $\mat{X}$ gives a discrete approximation of the
parameterized free-surface manifold
\[
    \mathcal{M}_h
    =
    \{\,\mat{h}(\params): \params \in \mathcal{D}\,\}.
\]

\begin{algorithm}[!ht]
\caption{POD basis construction.}
\label{alg:SVD}
\begin{algorithmic}[1]
\State \textbf{Input:} centered snapshot matrix
$\tilde{\mat{X}} \in \mathbb{R}^{N_h \times M}$,
POD tolerance $\epsilon_{\mathrm{pod}}$.
\State Compute the economy singular value decomposition
\begin{equation*}
\tilde{\mat{X}} = \mat{U}\,\mat{\Sigma}\,\mat{V}^{\top},
\quad
\mat{U} \in \mathbb{R}^{N_h \times M},\;\;
\mat{V} \in \mathbb{R}^{M \times M},\;\;
\mat{\Sigma} = \mathrm{diag}(\sigma_1,\dots,\sigma_M),
\end{equation*}
with $\sigma_1 \ge \sigma_2 \ge \cdots \ge \sigma_M \ge 0$.
\State Normalize the singular-value spectrum:
$\bar{\sigma}_i = \sigma_i / \sigma_1$, $i = 1,\dots,M$.
\State Select the reduced dimension
\[
r =
\begin{cases}
\max\{\, i : \bar{\sigma}_i > \epsilon_{\mathrm{pod}} \,\} & \text{if } \bar{\sigma}_M \le \epsilon_{\mathrm{pod}}, \\
M & \text{otherwise.}
\end{cases}
\]
\State Form the reduced basis from the leading $r$ left singular vectors:
\[
\mat{\Phi}_r \;=\; \mat{U}(:,\,1\!:\!r) \;\in\; \mathbb{R}^{N_h \times r}.
\]
\State \textbf{Output:} reduced basis matrix $\mat{\Phi}_r$ and reduced dimension $r$.
\end{algorithmic}
\end{algorithm}
POD constructs a low-dimensional orthonormal basis that captures the
dominant structures of the free-surface snapshot data through the singular
value decomposition (SVD) of the snapshot matrix. Because $\mat{X}$
typically contains a dominant mean component together with a smaller
fluctuating component, a direct SVD would force the leading mode to
represent the mean free-surface profile. To focus the basis on the
parametric variations of the free surface, the mean snapshot
\[
\bar{\mat{h}} = \frac{1}{M}\sum_{i=1}^{M} \mat{h}(\params_i)
\]
is subtracted first, and the centered snapshot matrix
\begin{equation}
\tilde{\mat{X}} =
\bigl[\, \mat{h}(\params_1) - \bar{\mat{h}},\;\, \mat{h}(\params_2) - \bar{\mat{h}},\;\, \dots,\;\, \mat{h}(\params_M) - \bar{\mat{h}} \,\bigr]
\end{equation}
is decomposed as
\begin{equation}
\tilde{\mat{X}} = \mat{\Phi}\,\mat{\Sigma}\,\mat{V}^{\top},
\end{equation}
where $\mat{\Phi} \in \mathbb{R}^{N_h \times N_h}$ and $\mat{V} \in \mathbb{R}^{M \times M}$ are orthogonal, and $\mat{\Sigma}$ contains the singular values $\sigma_1 \ge \sigma_2 \ge \cdots \ge 0$.

 The truncation rank is determined from the normalized singular value spectrum through a user-prescribed tolerance $\epsilon_{\mathrm{pod}}$: $r$ is taken as the largest index for which the normalized singular value $\sigma_i / \sigma_1$ exceeds $\epsilon_{\mathrm{pod}}$. The procedure is summarized in Algorithm~\ref{alg:SVD}.


Once the POD basis $\mat{\Phi}_r \in \mathbb{R}^{N_h \times r}$ is fixed, the
reduced coordinates $\mat{a}(\newparams)$ for a new parameter
$\newparams \in \mathcal{P}$ would still require a full-order solution under a
Galerkin-projection approach. To eliminate this cost, a Gaussian Process
Regression (GPR) model $\hat{\mat{\pi}} : \mathbb{R}^d \to \mathbb{R}^{r}$ is
trained offline on the dataset
$\mathcal{D} = \{(\params_i,\, \mat{a}(\params_i)) : i = 1, \dots, M\}$ to learn
the parameter-to-coefficient mapping
\begin{equation}
  \params \mapsto \mat{a}(\params)
  = \mat{\Phi}_r^\top \bigl(\mat{h}(\params) - \bar{\mat{h}}\bigr)
  \approx \hat{\mat{a}}(\params).
\end{equation}
Each of the $r$ reduced coordinates is modelled as an independent Gaussian
process, placing a prior over functions and conditioning on the training data to
yield a posterior predictive mean $\hat{\mat{a}}(\newparams)$ (together with a
variance quantifying predictive uncertainty).

Once trained, the ROM prediction of the free-surface profile for any new
$\newparams \in \mathcal{P}$ follows entirely from the surrogate as
\begin{equation}
  \mat{h}_{\mathrm{rom}}(\newparams)
  = \bar{\mat{h}}
  + \mat{\Phi}_r\,\hat{\mat{a}}(\newparams),
  \label{eq:rom_prediction}
\end{equation}
bypassing the full-order model during online evaluation.

% ==============================================================================
% REVISED SECTION: inference procedure with the error model
% (model discrepancy, mesh discretization error; ROM error verified negligible)
% ==============================================================================
\section{Inference procedure}
\label{sec:bayesian_inference}

The forward model of Section~\ref{sec:rom} defines the deterministic map from
the physical parameter vector $\params$ to the predicted free-surface profile,
$\vec{d}(\params) = \mat{h}_{\mathrm{rom}}(\params)$, evaluated with $r$ POD
modes according to Eq.~\eqref{eq:rom_prediction}. For an Oldroyd-B fluid,
$\params=(\lambda,\,\beta)^{\top}$, while for a Giesekus fluid,
$\params=(\lambda,\,\beta,\,\alpha)^{\top}$. The general Bayesian
framework, including the observation model~\eqref{eq:error_decomposition} and
Bayes' theorem~\eqref{eq:bayes_intro}, was introduced in
Section~\ref{sec:bayes}. We now specify the error model
(Section~\ref{subsec:error_model}), the priors (Section~\ref{subsec:priors}),
the likelihood (Section~\ref{subsec:likelihood}), and the MCMC sampler
(Section~\ref{subsec:posterior}).

\subsection{Error model}
\label{subsec:error_model}

Of the four error terms of the observation
model~\eqref{eq:error_decomposition}, the ROM error is shown below to be
negligible and is omitted from the inference. The three remaining terms are
modeled as zero-mean Gaussian random vectors that are independent of each other,
\begin{equation}
\vec{\varepsilon}_{\subexp}
\sim \mathcal{N}\!\left(\vec{0},\,\sigma_{\subexp}^{2}\,\mat{I}\right),
\qquad
\vec{\varepsilon}_{\subbias}
\sim \mathcal{N}\!\left(\vec{0},\,\sigma_{\subbias}^{2}\,\mat{K}'\right),
\qquad
\vec{\varepsilon}_{\submesh}
\sim \mathcal{N}\!\left(\vec{0},\,\mat{C}_{\submesh}\right).
\label{eq:exp_noise_prior}
\end{equation}
Here $\sigma_{\subexp}>0$ is the measurement-noise standard deviation,
$\sigma_{\subbias}>0$ is the amplitude of the correlated discrepancy with
correlation matrix $\mat{K}'$, and $\mat{C}_{\submesh}$ is the covariance matrix
of the mesh discretization error. The terms are summarized in
Table~\ref{tab:error_terms}; the measurement noise follows
Eq.~\eqref{eq:exp_noise_prior} directly, and the remaining terms are specified
below.

\begin{table}[!ht]
\centering
\small
\begin{tabular}{llll}
\toprule
Term & Distribution & Origin & Inferred \\
\midrule
$\vec{d}(\params)$
  & deterministic
  & ROM prediction with $r$ POD modes
  & $\params$ \\
$\vec{\varepsilon}_{\rom}(\params)$
  & neglected
  & POD truncation and GPR error
  & no \\
$\vec{\varepsilon}_{\submesh}$
  & $\mathcal{N}\bigl(\vec{0},\,\mat{C}_{\submesh}\bigr)$
  & mesh discretization of the FOM
  & no \\
$\vec{\varepsilon}_{\subbias}$
  & $\mathcal{N}\bigl(\vec{0},\,\sigma_{\subbias}^{2}\mat{K}'\bigr)$
  & constitutive-model discrepancy
  & $\sigma_{\subbias}$ \\
$\vec{\varepsilon}_{\subexp}$
  & $\mathcal{N}\bigl(\vec{0},\,\sigma_{\subexp}^{2}\mat{I}\bigr)$
  & measurement noise
  & $\sigma_{\subexp}$ \\
\bottomrule
\end{tabular}
\caption{Terms of the observation model~\eqref{eq:error_decomposition}.}
\label{tab:error_terms}
\end{table}

\paragraph{Model discrepancy}
The discrepancy is modeled as a zero-mean Gaussian process along the extrudate,
$\varepsilon_{\subbias} \sim \mathcal{GP}\bigl(0,\,\sigma_{\subbias}^{2}\,
k_{\subbias}\bigr)$, with the squared-exponential correlation
function~\cite{rasmussen2006gpml}
\begin{equation}
k_{\subbias}(x,x')
= \exp\!\left(-\dfrac{(x - x')^{2}}{2\,\ell_{\subbias}^{2}}\right),
\qquad
\mat{K}_{ij} = k_{\subbias}(x_{i},x_{j}).
\label{eq:kernel}
\end{equation}
The parameter $\ell_{\subbias}$ is the correlation length of the discrepancy,
in the same non-dimensional units as $x_i$. An unconstrained discrepancy is
confounded with the physical parameters~\cite{brynjarsdottir2014}. To reduce
this confounding, the process is conditioned on two physical constraints: the
free surface is pinned at the die exit, $\varepsilon_{\subbias}(0)=0$, and the
extrudate is flat far downstream, $\varepsilon_{\subbias}'(x^{c}_{j})=0$ at
$N_c$ points $x^{c}_{j}\in[3.5,\,5]$. The squared-exponential kernel is
differentiable, so that the derivative of the process is again a Gaussian
process, with
\begin{equation}
\begin{aligned}
\operatorname{Cov}\bigl(\varepsilon_{\subbias}(x),\,
                        \varepsilon_{\subbias}'(x')\bigr)
&= \sigma_{\subbias}^{2}\,\frac{x-x'}{\ell_{\subbias}^{2}}\,
   k_{\subbias}(x,x'),
\\
\operatorname{Cov}\bigl(\varepsilon_{\subbias}'(x),\,
                        \varepsilon_{\subbias}'(x')\bigr)
&= \sigma_{\subbias}^{2}\,\frac{1}{\ell_{\subbias}^{2}}
   \left(1-\frac{(x-x')^{2}}{\ell_{\subbias}^{2}}\right)
   k_{\subbias}(x,x').
\end{aligned}
\label{eq:kernel_derivatives}
\end{equation}
Conditioning on the $N_c+1$ constraints gives the correlation matrix at the
observation points,
\begin{equation}
\mat{K}'
= \mat{K} - \mat{K}_{c}\,\mat{K}_{cc}^{-1}\,\mat{K}_{c}^{\top},
\label{eq:kernel_constrained}
\end{equation}
where $\mat{K}_{c}\in\mathbb{R}^{N_h\times(N_c+1)}$ holds the correlations
between $\varepsilon_{\subbias}(x_i)$ and the constrained quantities, and
$\mat{K}_{cc}\in\mathbb{R}^{(N_c+1)\times(N_c+1)}$ those among the constrained
quantities, both evaluated from Eqs.~\eqref{eq:kernel}
and~\eqref{eq:kernel_derivatives} with $\sigma_{\subbias}=1$.

\paragraph{Mesh discretization error}
A real measurement carries no discretization error: all of it resides in the
forward model. The synthetic experiments are set up to reproduce this
situation. The observation is computed on the reference mesh
$m_{\mathrm{ref}}$, the finest mesh of the convergence study, whereas the
snapshots used to train the ROM are computed on the coarser mesh $m_2$
($\Delta x_{\mathrm{box}} = 0.1$, $\Delta x_{\mathrm{wall}} = 0.02$). The
discretization error of $m_2$ therefore enters the residual in full, instead of
cancelling as it would if the observation and the snapshots shared the same
mesh. The observed profile is interpolated to the stations $x_i$ of the mesh
$m_2$. The mesh error is modeled as a zero-mean Gaussian process whose amplitude
follows the swell displacement $w(x) = h_{\obs}(x) - 1$,
\begin{equation}
k_{\submesh}(x,x')
= \sigma_{\submesh}^{2}\,w(x)\,w(x')
  \exp\!\left(-\dfrac{(x - x')^{2}}{2\,\ell_{\submesh}^{2}}\right),
\qquad
\bigl(\mat{C}_{\submesh}\bigr)_{ij} = k_{\submesh}(x_{i},x_{j}),
\label{eq:kernel_mesh}
\end{equation}
with the relative amplitude $\sigma_{\submesh}$ and the correlation length
$\ell_{\submesh}$. Since $w(0)=0$, the mesh error vanishes at the die exit
without any conditioning. The hyperparameters $\sigma_{\submesh}$ and
$\ell_{\submesh}$ are not inferred from $\vec{h}_{\obs}$, because they would be
confounded with the model discrepancy; they are fitted to a mesh convergence
study instead. For each case $q=1,\dots,Q$ of the study (values of $\beta$ and
$\wi$, distinct from the cases used for the inference), the error curve of the
mesh $m_2$ is
\begin{equation}
\vec{e}_{\submesh,q} = \vec{h}_{m_2,q} - \vec{h}_{\mathrm{ref},q},
\label{eq:mesh_error_curves}
\end{equation}
where $\vec{h}_{\mathrm{ref},q}$ is the free surface computed on the reference
mesh, and all profiles are interpolated to the observation points. The
hyperparameters then follow from the maximum-likelihood estimate
\begin{equation}
\bigl(\hat{\sigma}_{\submesh},\,\hat{\ell}_{\submesh}\bigr)
= \arg\max_{\sigma_{\submesh},\,\ell_{\submesh}}\;
  \sum_{q=1}^{Q}
  \log \mathcal{N}\Bigl(\vec{e}_{\submesh,q} \,\Big|\, \vec{0},\;
  \mat{C}_{\submesh}\bigl(\sigma_{\submesh},\ell_{\submesh};\vec{w}_{q}\bigr)
  + \sigma_{\mathrm{ref}}^{2}\,\mat{I}\Bigr),
\label{eq:mesh_mle}
\end{equation}
where $\vec{w}_{q} = \vec{h}_{\mathrm{ref},q} - 1$ and the nugget
$\sigma_{\mathrm{ref}}$ represents the error of the reference mesh itself. It
is taken from the difference between the two finest meshes; since the grid
convergence index (GCI) of~\cite{celik2008procedure} is approximately a 95\%
bound, $\sigma_{\mathrm{ref}}\approx\mathrm{GCI}/2$. Because the observation is
computed on $m_{\mathrm{ref}}$, the error curves $\vec{e}_{\submesh,q}$ are, up
to their sign, the contribution that enters the
residual~\eqref{eq:residual}.

In the limit $\ell_{\submesh}\to\infty$ the mesh covariance has rank one,
$\mat{C}_{\submesh} = \sigma_{\submesh}^{2}\,\vec{w}\vec{w}^{\top}$, and the
mesh error reduces to a rescaling of the swell displacement,
$\vec{h}_{\obs} - 1 = (1+\kappa)\,(\vec{d}(\params) - 1) + \dots$ with
$\kappa\sim\mathcal{N}(0,\sigma_{\submesh}^{2})$. An alternative to a fixed
$\mat{C}_{\submesh}$ is to treat the mesh size $\Delta$ as an additional input
of the emulator, $h(x,\params,\Delta) = h_{\infty}(x,\params)
+ e(x,\params,\Delta)$, with a Gaussian process for $e$ whose covariance
vanishes as $\Delta\to0$, \eg $k(\Delta,\Delta')\propto(\Delta\Delta')^{p}$,
trained on coarse runs over the whole parameter range and a few fine
runs~\cite{kennedy2000predicting,tuo2014surrogate}. The emulator then predicts
the converged surface $h_{\infty}$, and its predictive variance includes the
discretization uncertainty at every $\params$. This extension is not pursued
here.

\paragraph{Reduced-order model error}
The ROM error has two parts: a truncation part, which decreases with the number
of retained modes $r$, and a regression part, which stems from the GPR
prediction of the reduced coefficients. Both are measured together by the
leave-one-out errors
\begin{equation}
\vec{e}_{\rom,i}
= \mat{h}(\params_i) - \mat{h}_{\mathrm{rom}}^{(-i)}(\params_i),
\qquad i=1,\dots,M,
\label{eq:loo_errors}
\end{equation}
where $\mat{h}_{\mathrm{rom}}^{(-i)}$ is the ROM with $r$ modes trained without
the snapshot $i$. Each snapshot is thus predicted by a ROM that has not seen
it, which makes these errors a conservative estimate of the ROM error at a new
parameter value. Their root-mean-square value,
\begin{equation}
\mathrm{RMS}_{\rom}
= \left(\frac{1}{M N_h}\sum_{i=1}^{M}
  \bigl\|\vec{e}_{\rom,i}\bigr\|^{2}\right)^{1/2},
\label{eq:loo_rms}
\end{equation}
amounts to \rai{X} (\rai{Y}\% of the maximum swell displacement). This is
\rai{more than an order of magnitude} smaller than the measurement noise
(\rai{value} at the lowest noise level considered), the mesh error
(\rai{value}) and the model discrepancy (\rai{value}). The ROM error is
therefore neglected, $\vec{\varepsilon}_{\rom}=\vec{0}$, and the ROM prediction
$\vec{d}(\params)$ is treated as the full-order solution on the mesh $m_2$.

\subsection{Prior distributions}
\label{subsec:priors}

The prior encodes what is known before the profile is observed. For the
physical parameters, we assume positivity and broad, physically reasonable
ranges. We place a uniform prior on these variables, bounded by the range over
which the reduced-order model has been trained.

The remaining unknowns are the amplitudes of the measurement noise and of the
model discrepancy. Their priors are
\begin{equation}
\sigma_{\subexp}\sim \mathcal{E}(\gamma_{\subexp}),
\qquad
\sigma_{\subbias}\sim \mathcal{E}(\gamma_{\subbias}),
\label{eq:priors_noise}
\end{equation}
where $\gamma_{\subexp}$ and $\gamma_{\subbias}$ are the rate parameters of the
exponential distributions. These priors favor small amplitudes while still
allowing larger values when supported by the data. The prior of
$\sigma_{\subbias}$ is centred on the root-mean-square value
$\hat{\sigma}_{\subbias} = \mathrm{RMS}(\hat{\vec{\varepsilon}}_{\subbias})$ of
the residual $\hat{\vec{\varepsilon}}_{\subbias}$ of a maximum-likelihood fit
of the ROM to the noise-free observation, \ie
$\gamma_{\subbias} = 1/\hat{\sigma}_{\subbias}$. This fit includes
$\mat{C}_{\submesh}$ in its covariance, so that $\sigma_{\subbias}$ does not
absorb the mesh error a second time.

When model bias is inferred, $\ell_{\subbias}$ is fixed and
$\sigma_{\subexp}$ and $\sigma_{\subbias}$ are inferred jointly with
$\params$. When the data are generated with the same constitutive model as the
forward model, the discrepancy term is omitted and only $\sigma_{\subexp}$ is
inferred alongside $\params$. The covariance $\mat{C}_{\submesh}$ carries no
prior, since its hyperparameters are fixed before the sampling
(Section~\ref{subsec:error_model}). Case-specific numerical values
are reported in Section~\ref{sec:numerical_experiments}.

\subsection{Likelihood function}
\label{subsec:likelihood}

Because all error terms are Gaussian and independent, they can be integrated
out analytically: their sum
$\vec{\varepsilon}=\vec{\varepsilon}_{\subexp}+\vec{\varepsilon}_{\subbias}
+\vec{\varepsilon}_{\submesh}$ is again Gaussian, with the covariance
\begin{equation}
\mat{C}(\sigma_{\subexp},\sigma_{\subbias})
= \sigma_{\subexp}^{2}\,\mat{I}
+ \sigma_{\subbias}^{2}\,\mat{K}'
+ \mat{C}_{\submesh}
\,\in\, \mathbb{R}^{N_h\times N_h},
\label{eq:covariance}
\end{equation}
where $\mat{I}$ is the identity matrix, $\mat{K}'$ is the constrained
correlation matrix from Eq.~\eqref{eq:kernel_constrained}, and
$\mat{C}_{\submesh}$ follows from Eq.~\eqref{eq:kernel_mesh}. All matrices have
size $N_h\times N_h$ for a profile with $N_h$ observations.

The residual is given by
\begin{equation}
\vec{r}(\params)
= \vec{h}_{\obs} - \vec{d}(\params),
\qquad
\vec{r},\,\vec{h}_{\obs},\,\vec{d}\in \mathbb{R}^{N_h}.
\label{eq:residual}
\end{equation}
The likelihood is the density of the combined error evaluated at this residual:
\begin{equation}
\mathcal{L}\bigl({\params},\,\sigma_{\subexp},\,\sigma_{\subbias}
\mid \vec{h}_{\obs}\bigr)
= \dfrac{1}{(2\pi)^{N_h/2}\,|\mat{C}|^{1/2}}
  \exp\!\left(-\tfrac{1}{2}\,\vec{r}^{\top}\mat{C}^{-1}\vec{r}\right).
\label{eq:likelihood}
\end{equation}
The corresponding log-likelihood is
\begin{equation}
\log \mathcal{L}\bigl({\params},\,\sigma_{\subexp},\,\sigma_{\subbias}
\mid \vec{h}_{\obs}\bigr)
= -\tfrac{1}{2}\,\vec{r}^{\top}\mat{C}^{-1}\vec{r}
  -\tfrac{1}{2}\,\log|\mat{C}|
  -\tfrac{N_h}{2}\log(2\pi).
\label{eq:loglikelihood}
\end{equation}
It is evaluated with the Cholesky factorization
$\mat{C}=\mat{G}\mat{G}^{\top}$, with $\mat{G}$ lower triangular, from which
$\vec{r}^{\top}\mat{C}^{-1}\vec{r} = \|\mat{G}^{-1}\vec{r}\|^{2}$ and
$\log|\mat{C}| = 2\sum_{i}\log \mat{G}_{ii}$.

In the rank-one limit of the mesh error,
$\mat{C}_{\submesh} = \sigma_{\submesh}^{2}\,\vec{w}\vec{w}^{\top}$, the mesh
contribution can be separated from the remaining covariance
$\mat{C}_{0} = \mat{C} - \mat{C}_{\submesh}$ with the Woodbury identity,
\begin{equation}
\vec{r}^{\top}\mat{C}^{-1}\vec{r}
= \vec{r}^{\top}\mat{C}_{0}^{-1}\vec{r}
- \frac{\sigma_{\submesh}^{2}\,
        \bigl(\vec{w}^{\top}\mat{C}_{0}^{-1}\vec{r}\bigr)^{2}}
       {1+\sigma_{\submesh}^{2}\,\vec{w}^{\top}\mat{C}_{0}^{-1}\vec{w}},
\qquad
\log|\mat{C}|
= \log|\mat{C}_{0}|
+ \log\bigl(1+\sigma_{\submesh}^{2}\,
            \vec{w}^{\top}\mat{C}_{0}^{-1}\vec{w}\bigr).
\label{eq:woodbury}
\end{equation}
The remaining covariance
$\mat{C}_{0} = \sigma_{\subexp}^{2}\,\mat{I} + \sigma_{\subbias}^{2}\,\mat{K}'$
is diagonalized by the eigenvectors of $\mat{K}'$, which are computed once
before the sampling.

\subsection{Posterior sampling}
\label{subsec:posterior}

Bayes' theorem~\eqref{eq:bayes_intro} specifies the posterior. The inferred
quantities are the material parameters $\params$ and the error hyperparameters
$\sigma_{\subexp}$ and $\sigma_{\subbias}$; $\mat{C}_{\submesh}$ is fixed
before the sampling. With independent priors,
\begin{equation}
p\bigl({\params},\sigma_{\subexp},\sigma_{\subbias}
      \mid \vec{h}_{\obs}\bigr)
\;\propto\;
\mathcal{L}\bigl({\params},\,\sigma_{\subexp},\,\sigma_{\subbias}
\mid \vec{h}_{\obs}\bigr)\,
p(\params)\,p(\sigma_{\subexp})\,p(\sigma_{\subbias}).
\label{eq:posterior_full}
\end{equation}
The evidence $p(\vec{h}_{\obs})$ only adds a constant, so the target
density is
\begin{equation}
\log p\bigl({\params},\sigma_{\subexp},\sigma_{\subbias}
      \mid \vec{h}_{\obs}\bigr)
= \log \mathcal{L}\bigl({\params},\,\sigma_{\subexp},\,\sigma_{\subbias}
\mid \vec{h}_{\obs}\bigr)
+ \log p\bigl({\params},\sigma_{\subexp},\sigma_{\subbias}\bigr)
+ \mathrm{const},
\label{eq:log_posterior}
\end{equation}
where the constant $-\log p(\vec{h}_{\obs})$ is omitted in practice. For the
Oldroyd-B model, the sampling can equivalently be carried out in the
parametrization $(\theta_1,\theta_2) = ((1-\beta)\lambda,\,\beta)$, in which
$\theta_1$ is the combination that sets the stress ratio $N_1/\tau_w$. Because
the forward models are nonlinear and the covariance contains inferred
hyperparameters, this posterior has no closed form. We therefore sample it
using Markov chain Monte Carlo. Specifically, we use the Affine Invariant
Stretch Move sampler of Goodman and Weare~\cite{goodman2010ensemble},
implemented in \texttt{emcee}~\cite{foremanmackey2013emcee}. Multiple walkers
explore the posterior in parallel, which helps reduce sensitivity to their
starting positions. Further implementation details follow
Ref.~\cite{rinkens2023uncertainty}.

In all experiments, we use $N_{\mathrm{w}}=2N_{\mathrm{dim}}$ walkers and
$10{,}000$ iterations per walker. The first $30\%$ of samples are discarded as
burn-in, allowing the walkers to move from their prior initialization to the
typical set. Convergence is assessed using trace plots and the Gelman-Rubin
statistic~\cite{gelman1992rubin,brooks1998convergence}. We accept a chain when
the walkers mix around a common stationary region, $\hat{R}$ is close to one,
and in practice $\hat{R}\lesssim1.1$, and the posterior means, credible
intervals, and correlations remain stable as the chain length is increased.
% ==============================================================================
% end of revised section
% ==============================================================================


\FloatBarrier
\section{Numerical experiments}
\label{sec:numerical_experiments}

\subsection{ROM performance}
\begin{itemize}
\item the value of $\epsilon_{\text{pod}}$, the retained rank $r$ and singular value spectrum
\item Number of points along the profile
\item offline cost (450 Giesekus FOM runs) versus online cost, \ie the speedup that motivates the whole paper.
\item  ROM and FOM error comparison (for Giesekus, Oldoryd-B)

\end{itemize}

\begin{table}[!ht]
  \centering
  \caption{Planned parameter spaces for the three constitutive models.}
  \label{tab:parameter-spaces}
  \begin{tabular}{llll}
    \toprule
    Model & Inferred parameters & Training columns & Example bounds \\
    \midrule
    Oldroyd-B & $\lambda,\beta$ & 2 & $\lambda\in[1,10]$, $\beta\in[0.1,0.9]$ \\
    Giesekus & $\lambda,\beta,\alpha$ & 3 & add $\alpha\in[0.1,0.5]$ \\
    PTT & $\lambda,\beta,\epsilon$ & 3 & add $\epsilon\in[\epsilon_{\min},\epsilon_{\max}]$ \\
    \bottomrule
  \end{tabular}
\end{table}
\subsection{Tanner expression}

For an axisymmetric extrudate issuing from a circular die, Tanner's
die-swell theory relates the final swell ratio to the elastic recovery at
the die wall \cite{Tanner1970}. If \(B\) denotes the measured swell ratio,
Tanner's expression can be written as
\begin{equation}
    B
    =
    0.13
    +
    \left[
        1
        +
        \frac{1}{2}
        \left(
            \frac{N_{1}}{2\tau_w}
        \right)^2
    \right]^{1/6},
    \label{eq:tanner_swell}
\end{equation}
where \(N_{1}\) is the first normal-stress difference at the die wall and
\(\tau_w\) is the corresponding wall shear stress. The constant
\(0.13\) is a correction added to account for the
die swell observed for a Newtonian fluid, so that in the limit
\(N_{1}\rightarrow 0\) the expression predicts \(B \approx 1.13\), rather
than unity.

The relation is applicable in the low-Weissenberg-number regime, where the
assumptions underlying Tanner's approximation remain valid \cite{comminal2018numerical}. Following \cite{comminal2018numerical}, there is an analytical expression relating $N_1/(2 \tau_w) = (1-\beta) \lambda \dot{\gamma}_{\text{w}}$ for the fully-developed stress profiles for the Oldroyd-B model where $\dot{\gamma}_{\text{w}}$ is the wall shear rate.  Furthermore, a measurement of
the final swell ratio provides information only about the dimensionless
stress ratio \(N_{1}/\tau_w\), rather than \(N_{1}\) and \(\tau_w\)
independently.

\subsection{Inference cases}

\emph{Case I: Oldroyd-B}


\begin{table}[!ht]
\centering
\begin{tabular}{ll}
\toprule
Parameter & Prior \\
\midrule
${\lambda}$ & $\mathcal{U}(1,\,10)$ \\
$\beta$ & $\mathcal{U}(0.1,\,0.9)$ \\
$\sigma_{\mathrm{exp}}$      & $\mathcal{E}(8)$ \\
\bottomrule
\end{tabular}
\caption{Prior distributions used for the Oldroyd-B inference cases.}
\label{tab:prior_oldoryd}
\end{table}

\begin{itemize}
\item generated 90 snapshots (h(x) curves) with 10 x 9 linearly spaced lambda values between [1, 10] and beta values between [0.1, 0.9] with $U_\text{avg} = 0.2$.
    \item trained the ROM with the snapshots with GPR.

\begin{figure}[!ht]
\centering

\begin{subfigure}{0.32\linewidth}
    \centering
    \includegraphics[width=\linewidth]{plots/oldroyd/noise-0.5/rom_prediction_noise_0.5.pdf}
    \caption{Noise = 0.5\%}
\end{subfigure}
\begin{subfigure}{0.32\linewidth}
    \centering
    \includegraphics[width=\linewidth]{plots/oldroyd/noise-2/rom_prediction_noise_2.pdf}
    \caption{Noise = 2\%}
\end{subfigure}
\begin{subfigure}{0.32\linewidth}
    \centering
    \includegraphics[width=\linewidth]{plots/oldroyd/noise-5/rom_prediction_noise_5.pdf}
    \caption{Noise = 5\%}
\end{subfigure}

\caption{ROM comparison with the FOM data for the Oldroyd model with synthetic noise prescribed as a percentage of the maximum swelling displacement.}
\label{fig:rom_fom_oldroyd_noise}
\end{figure}

    \item generated the fully swollen steady state curve h(x) using FOM for the Oldroyd B model for $\lambda = 5.5$ and $\beta = 0.65$ with $U_\text{avg} = 0.2$ to infer the parameter.
    \item generated 56 equally spaced points used for the inference.
    \item added noise to the curve h(x) by adding 2\% of the max displacement to mimic the synthetic data.
    \item Compared the fully swollen FOM for $\lambda = 5.5$ and $\beta = 0.65$ data with the ROM and the comparison looks good in Figure~\ref{fig:rom_fom_oldroyd_noise}.
    \item The prior for $\sigma_{\text{exp}}$ is obtained so that the mean of the exponential = 10\% of the maximum swell height and the uniform prior for the $(\lambda, \beta)$ are used within the training regimes for the reduced model. The prior is listed in Table~\ref{tab:prior_oldoryd}.
    \item Now we infer ($\lambda$, $\beta$) and the corner plot shown in Figure~\ref{fig:corner_oldrodyd} shows the prior is updated to get a good posterior, which is centered close to ground truth. There is a positive correlation between $\lambda$ and $\beta$ which can be explained by $N_1/(2 \tau_{\text{w}}) = (1 -\beta) \lambda \dot{\gamma}_{\text{w}}$.

    \item We also tried inferring the parameter ($\lambda$, $\beta$) using only the swell height ratio, but the posterior remains the same as the prior, meaning that there is not much information in the likelihood to update the posterior.

    \item We also look at the noise sensitivity with $\sigma_{\text{noise}} = \{0.1, 0.5, 1, 2, 5, 10\}$ percent of the maximum swell height.

\begin{table}[htbp]
\centering
\begin{tabular}{lcccccc}
\toprule
Noise &
mean($\lambda$) & std($\lambda$) &
mean($\beta$) & std($\beta$) &
mean($\sigma_{\mathrm{noise}}$) & std($\sigma_{\mathrm{noise}}$) \\
\midrule
0.1\%  & 5.50 & $1.29 \times 10^{-2}$ & 0.648 & $1.72 \times 10^{-3}$ & $7.50 \times 10^{-4}$ & $1.26 \times 10^{-4}$ \\
0.5\%  & 5.49 & $5.55 \times 10^{-2}$ & 0.646 & $7.40 \times 10^{-3}$ & $3.28 \times 10^{-3}$ & $5.32 \times 10^{-4}$ \\
1\%    & 5.46 & $1.09 \times 10^{-1}$ & 0.644 & $1.47 \times 10^{-2}$ & $6.42 \times 10^{-3}$ & $1.02 \times 10^{-3}$ \\
2\%    & 5.42 & $2.08 \times 10^{-1}$ & 0.638 & $2.87 \times 10^{-2}$ & $1.26 \times 10^{-2}$ & $1.94 \times 10^{-3}$ \\
5\%    & 5.23 & $4.97 \times 10^{-1}$ & 0.605 & $8.19 \times 10^{-2}$ & $3.10 \times 10^{-2}$ & $4.74 \times 10^{-3}$ \\
10\%   & 4.89 & $7.66 \times 10^{-1}$ & 0.514 & $1.69 \times 10^{-1}$ & $6.12 \times 10^{-2}$ & $9.16 \times 10^{-3}$ \\
\bottomrule
\end{tabular}
\caption{Effect of the prescribed noise percentage on the inferred parameters for the Oldroyd model. The true parameter values listed in the source file are $\lambda = 5.5$ and $\beta = 0.65$.}
\label{tab:oldroyd-noise-sensitivity}
\end{table}

\end{itemize}



\begin{figure}[!ht]
    \centering
    \includegraphics[width=0.7\linewidth]{plots/oldroyd/noise-2/corner_physical_all_noise_2.pdf}
    \caption{corner plots}
    \label{fig:corner_oldrodyd}
\end{figure}

\FloatBarrier

\emph{Case II: Giesekus}

\begin{table}[!ht]
\centering
\begin{tabular}{ll}
\toprule
Parameter & Prior \\
\midrule
${\lambda}$ & $\mathcal{U}(1,\,10)$ \\
$\gamma$ & $\mathcal{U}(0.1,\,0.9)$ \\
$\alpha$ & $\mathcal{U}(0.1,\,0.5)$ \\
$\sigma_{\mathrm{exp}}$
& $\mathcal{E}(8)$ \\
\bottomrule
\end{tabular}
\caption{Prior distributions used for the Giesekus inference cases.}
\label{tab:prior_giesekus}
\end{table}


\begin{figure}[!ht]
\centering

\begin{subfigure}{0.32\linewidth}
    \centering
    \includegraphics[width=\linewidth]{plots/giesekus/noise-0.5/rom_prediction_noise_0.5.pdf}
    \caption{Noise = 0.5\%}
\end{subfigure}
\hfill
\begin{subfigure}{0.32\linewidth}
    \centering
    \includegraphics[width=\linewidth]{plots/giesekus/noise-2/rom_prediction_noise_2.pdf}
    \caption{Noise = 2\%}
\end{subfigure}
\hfill
\begin{subfigure}{0.32\linewidth}
    \centering
    \includegraphics[width=\linewidth]{plots/giesekus/noise-5/rom_prediction_noise_5.pdf}
    \caption{Noise = 5\%}
\end{subfigure}

\caption{ROM comparison with the FOM data for the Giesekus model with synthetic noise prescribed as a percentage of the maximum swelling displacement.}
\label{fig:rom_fom_giesekus_noise}
\end{figure}
\begin{itemize}
\item used 56 equally spaced points for the inference.
    \item generated the fully swollen steady-state curve h(x) using FOM for the Oldroyd B model for $\lambda = 5.5$ and $\beta = 0.65$, $\alpha = 0.25$ with $U_\text{avg} = 0.2$ to infer the parameter.
    \item added noise to the curve h(x) by adding 2\% of the max displacement to mimic the synthetic data.
    \item compared the fully swollen FOM for $\lambda = 5.5$ and $\beta = 0.65, \alpha = 0.25$ data with the ROM and the comparison looks good.
    \item Now we infer ($\lambda$, $\beta, \alpha$) and the corner plot shows the prior is updated to get a good posterior, which is centered close to ground truth. There is a positive correlation between $\lambda$ and $\beta$ which can be explained by $N_1/(2 \tau_{\text{w}}) = (1 -\beta) \lambda \dot{\gamma}_{\text{w}}$.

    \item We also look at the noise sensitivity with $\sigma_{\text{noise}} = \{0.1, 0.5, 1, 2, 5, 10\}$ percent of the maximum swell height.
\begin{table}[htbp]
\centering
\small
\begin{tabular}{lcccccccc}
\toprule
Noise &
mean($\lambda$) & std($\lambda$) &
mean($\beta$) & std($\beta$) &
mean($\alpha$) & std($\alpha$) &
mean($\sigma_{\mathrm{noise}}$) & std($\sigma_{\mathrm{noise}}$) \\
\midrule
0.1\%  & 5.45 & $8.63 \times 10^{-2}$ & 0.648 & $1.58 \times 10^{-3}$ & 0.245 & $4.53 \times 10^{-3}$ & $1.70 \times 10^{-4}$ & $1.62 \times 10^{-5}$ \\
0.5\%  & 5.30 & $3.55 \times 10^{-1}$ & 0.639 & $6.48 \times 10^{-3}$ & 0.244 & $1.87 \times 10^{-2}$ & $8.41 \times 10^{-4}$ & $7.82 \times 10^{-5}$ \\
1\%    & 5.27 & $6.96 \times 10^{-1}$ & 0.625 & $1.22 \times 10^{-2}$ & 0.250 & $3.75 \times 10^{-2}$ & $1.69 \times 10^{-3}$ & $1.60 \times 10^{-4}$ \\
2\%    & 5.31 & $1.19$                  & 0.594 & $2.07 \times 10^{-2}$ & 0.266 & $6.45 \times 10^{-2}$ & $3.40 \times 10^{-3}$ & $3.27 \times 10^{-4}$ \\
5\%    & 5.72 & $1.65$                  & 0.506 & $4.40 \times 10^{-2}$ & 0.315 & $9.15 \times 10^{-2}$ & $8.46 \times 10^{-3}$ & $8.19 \times 10^{-4}$ \\
10\%   & 5.91 & $1.68$                  & 0.388 & $8.22 \times 10^{-2}$ & 0.341 & $9.36 \times 10^{-2}$ & $1.68 \times 10^{-2}$ & $1.62 \times 10^{-3}$ \\
\bottomrule
\end{tabular}
\caption{Effect of the prescribed noise percentage on the inferred parameters for the Giesekus model. The true parameter values listed in the source file are $\lambda = 5.5$, $\beta = 0.65$, and $\alpha = 0.25$.}
\label{tab:giesekus-noise-sensitivity}
\end{table}

  \item repeat inference at one or
  more average velocities $U$ to test whether a single profile is sufficient or
  whether joint inference over multiple profiles is needed.

\begin{figure}[!ht]
    \centering
    \includegraphics[width=0.7\linewidth]{plots/giesekus/noise-1/corner_physical_all_noise_1.pdf}
    \caption{corner plots}
    \label{fig:corner_giesekus}
\end{figure}
\end{itemize}

\FloatBarrier
\subsection{Model discrepancy}

\begin{figure}[!ht]
    \centering
    \begin{subfigure}[b]{0.33\linewidth}
        \centering
        \includegraphics[width=\linewidth]{plots/model_discrepancy/n1_vs_wi.pdf}
        \caption{First normal stress difference $N_1$.}
        \label{fig:n1_vs_wi}
    \end{subfigure}
    \begin{subfigure}[b]{0.34\linewidth}
        \centering
        \includegraphics[width=\linewidth]{plots/model_discrepancy/swell_vs_wi_beta_0.5.pdf}
        \caption{Swell ratio $h/H$.}
        \label{fig:swell_vs_wi}
    \end{subfigure}
\caption{Effect of the Giesekus mobility $\alpha$ on (a) $N_1$ at the die-wall shear rate $\dot{\gamma}_w = 4U_{\mathrm{avg}}/R$ and (b) the swell ratio $h/H$, versus $Wi = \lambda\dot{\gamma}_w$ ($\beta = 0.5$; $\alpha = 0$ is Oldroyd-B). Shear thinning limits the growth of $N_1$ and hence of the swell.}
    \label{fig:n1_swell_vs_wi}
\end{figure}

\begin{figure}[t]
  \centering
  \begin{subfigure}[b]{0.36\textwidth}
    \centering
    \includegraphics[width=\textwidth]{plots/model_discrepancy/model_bias_curves.pdf}
    \caption{Model bias $\varepsilon_{\mathrm{bias}}(x)$}
    \label{fig:model_bias_curves}
  \end{subfigure}\hfill
  \begin{subfigure}[b]{0.33\textwidth}
    \centering
    \includegraphics[width=\textwidth]{plots/model_discrepancy/acf_fit.pdf}
    \caption{Autocorrelation fit}
    \label{fig:acf_fit}
  \end{subfigure}\hfill
  \begin{subfigure}[b]{0.31\textwidth}
    \centering
    \includegraphics[width=\textwidth]{plots/model_discrepancy/l_bias_hist.pdf}
    \caption{Correlation length $\ell_{\mathrm{bias}}$}
    \label{fig:l_bias_hist}
  \end{subfigure}
  \caption{(a) Oldroyd-B model bias $\varepsilon_{\mathrm{bias}}(x) = y_{\mathrm{true}}(x) - y(x;\boldsymbol{\theta}_{\mathrm{MLE}})$ for all 280 Giesekus cases, coloured by $\alpha$; it grows with $\alpha$ and peaks near the die exit. (b) Sample autocorrelation of a representative $\varepsilon_{\mathrm{bias}}(x)$ (dots) and the fitted kernel $e^{-r^2/2\ell_{\mathrm{bias}}^2}$ (dashed); $\ell_{\mathrm{bias}}$ is the lag at which it drops to $e^{-1/2}$. (c) $\ell_{\mathrm{bias}}$ over all cases (mean $0.51$, dashed).}
  \label{fig:model_bias}
\end{figure}

\begin{figure}[t]
  \centering
  \begin{subfigure}[b]{0.33\textwidth}
    \centering
    \includegraphics[width=\textwidth]{plots/model_discrepancy/gp_unconstrained.pdf}
    \caption{Unconstrained GP}
    \label{fig:gp_unconstrained}
  \end{subfigure}
  \begin{subfigure}[b]{0.33\textwidth}
    \centering
    \includegraphics[width=\textwidth]{plots/model_discrepancy/gp_constrained.pdf}
    \caption{Constrained GP}
    \label{fig:gp_constrained}
  \end{subfigure}
\caption{Ten realizations of the model-bias GP ($\sigma_{\mathrm{bias}}^2 = 1$, $\ell_{\mathrm{bias}} = 0.51$) with $\pm\sigma(x)$ band (grey). (b) Constraining $\varepsilon_{\mathrm{bias}}(0) = 0$ (dot) and $\varepsilon_{\mathrm{bias}}' = 0$ on $[3.5, 5]$ (red) pins the bias at the die exit and flattens it downstream.}
  \label{fig:gp_constraints}
\end{figure}


\begin{figure}[!ht]
  \centering
  \begin{subfigure}[b]{0.38\textwidth}
    \centering
    \includegraphics[width=\textwidth]{plots/model_discrepancy/posterior_predictive_noise_without_bias_2.pdf}
    \caption{Without model discrepancy}
    \label{fig:pp_without_bias}
  \end{subfigure}
  \begin{subfigure}[b]{0.38\textwidth}
    \centering
    \includegraphics[width=\textwidth]{plots/model_discrepancy/posterior_predictive_noise_model_bias_2.pdf}
    \caption{With model discrepancy}
    \label{fig:pp_with_bias}
  \end{subfigure}
  \caption{Posterior predictive swell profile $h(x)$ of the Oldroyd-B ROM for Giesekus
    data ($\lambda = 11$, $\beta = 0.5$, $\alpha = 0.2$, 2\% noise; dots). (a) Without a
    discrepancy term the band is narrow but misses the data near the die exit
    ($0.5 \lesssim x \lesssim 1.5$). (b) With the constrained discrepancy term
    $\varepsilon_{\mathrm{bias}}$ the band widens to cover the data over the whole
    extrudate.}
  \label{fig:pp_bias}
\end{figure}

\begin{figure}[!ht]
    \centering
    \begin{subfigure}[b]{0.33\linewidth}
        \centering
        \includegraphics[width=\linewidth]{plots/model_discrepancy/n1_true_vs_inferred_no_bias.pdf}
        \caption{Without model discrepancy.}
        \label{fig:n1_inferred_no_bias}
    \end{subfigure}
    \begin{subfigure}[b]{0.33\linewidth}
        \centering
        \includegraphics[width=\linewidth]{plots/model_discrepancy/n1_true_vs_inferred_bias.pdf}
        \caption{With model discrepancy.}
        \label{fig:n1_inferred_bias}
    \end{subfigure}
    \caption{True first normal stress difference $N_1$ of the Giesekus data (solid lines, filled circles) and $N_1$ inferred with the Oldroyd-B reduced-order model (dashed lines, open squares) as functions of $Wi = \lambda\dot{\gamma}_w$, for $\alpha = 0.01$--$0.2$ and $\beta = 0.5$ (1\% noise). (a) Without a model-discrepancy term, the inferred $N_1$ increasingly overestimates the true value at high $Wi$ as shear thinning ($\alpha$) increases, by up to about 26\% for $\alpha = 0.2$. (b) With the constrained model-discrepancy term, the mean error for $Wi \geq 1.2$ is at most about 6\% for every $\alpha$. At $Wi = 0.4$, where the swell is close to Newtonian, $N_1$ is underestimated in both cases for $\alpha \geq 0.05$.}
    \label{fig:n1_inferred}
\end{figure}

\FloatBarrier
% \subsection{Synthetic experimental images}

% Synthetic experimental images are generated using Blender. The geometry is
% constructed as a surface of revolution from the prescribed axisymmetric
% profile $(x_i,r_i)$, which serves as the ground truth. The profile is revolved
% about the flow axis using 128 azimuthal segments. A short opaque die is added
% upstream of the free surface.

% An orthographic camera is positioned with its optical axis perpendicular to
% the flow axis. Consequently, the image scale is spatially uniform and no
% perspective correction or empirical camera calibration is required. The
% pixel-to-length conversion is determined from the known die radius
% $R_{\mathrm{die}}$. The main parameters used for the synthetic image
% generation are summarized in Table~\ref{tab:blender}.

% \begin{table}[!ht]
%   \centering
%   \caption{Parameters used for synthetic image generation. Lengths are
%   normalized by the die radius $R_{\mathrm{die}}$.}
%   \label{tab:blender}
%   \begin{tabular}{ll}
%     \hline
%     Software & Blender 4.2 \\
%     Renderer & Eevee \\
%     Geometry & surface of revolution \\
%     Azimuthal resolution & 128 segments \\
%     Axial resolution & prescribed by input profile \\
%     Camera & orthographic, optical axis $\perp$ flow axis \\
%     Orthographic scale & $8R_{\mathrm{die}}$ \\
%     Image resolution & $1200\times900$ px \\
%     \hline
%   \end{tabular}
% \end{table}
% \begin{figure}
% [!ht]
% \centering \includegraphics[width=0.33\linewidth]{swell.png} \caption{Synthetic experimental image generated in Blender. The video frames mimic an extrudate-swell experiment. Image processing extracts the fully swollen free-surface profile.} \label{fig:blender} \end{figure}
% Figure~\ref{fig:blender} shows an example of the synthetic experimental image.
% Since experimental images are not presently available, Blender is used to
% generate controlled image-like observations that mimic the imaging stage of
% an extrudate-swell experiment. The prescribed free-surface profile used to
% construct the Blender geometry serves only as ground truth and is not provided
% to the image-processing procedure. Instead, the free-surface profile is
% reconstructed directly from the rendered image and subsequently used as the
% observational data for parameter inference. This provides a controlled setting in which the complete pipeline of image acquisition, free-surface extraction, and parameter inference can be developed and validated before laboratory images become available.


% \subsection{Profile extraction}

% \begin{figure}[!ht]
%     \centering
%     \includegraphics[width=0.33\linewidth]{profile_from_image.pdf}
%     \caption{Extracted image from the blender}
%     \label{fig:image_extraction}
% \end{figure}
% The color of the extrudate is different from that of the die and background. The die and background are approximately gray and thus have low color saturation, while the extrudate has substantially higher saturation. The largest connected region in the image is considered to be the extrudate swell.

% The free-surface profile is extracted row by row in the direction
% transverse to the flow. For each axial image location $j$, the segmented
% extrudate intersects the row at the left and right boundaries,
% $u_j^{L}$ and $u_j^{R}$, respectively. Assuming an axisymmetric profile, the
% local radius in pixel units is obtained from half of the detected width.
% \begin{equation}
%   \rho_j = \frac{1}{2}\left(u_j^{R}-u_j^{L}\right).
% \end{equation}

% The known die radius is used to translate the pixel coordinates into non-dimensional physical coordinates. Let $a$ denote the number of
% pixels corresponding to one die radius, and let $j_0$ denote the axial pixel
% location of the die exit. The extracted coordinates are then written as
% \begin{equation}
%   x_j = \frac{j-j_0}{a},
%   \qquad
%   r_j = \frac{\rho_j}{a},
% \end{equation}
% such that both $x_j$ and $r_j$ are expressed in units of the die radius.
% Finally, the extracted radius profile is interpolated onto the prescribed
% observation grid $\{x_i\}$ used for the subsequent parameter inference.

\section{Conclusion}

The proposed method transforms extrudate swell from a forward simulation output
into an inverse material-characterization observable. The key computational
step is replacing the expensive FOM with a ROM that maps material parameters
directly to the observed free-surface curve. This makes MCMC feasible while
preserving the FOM physics through the offline snapshot database. The approach
is also compatible with experimental imaging because the likelihood only
requires a profile sampled on a known physical coordinate grid.

The model-family distinction is important. Oldroyd-B is a two-parameter inverse
problem and should be the first diagnostic case because it tests whether
$\lambda$ and $\beta$ can be recovered from the chosen profile. Giesekus and PTT
are three-parameter problems and are more likely to show parameter correlation
or non-identifiability. Comparing posterior predictive bands across model
families can reveal whether the data require the additional nonlinear
parameter, or whether Oldroyd-B already explains the measured swell profile
within the noise level.

The Blender workflow provides an intermediate validation layer between purely
synthetic FOM data and real experiments. Because the geometry, lighting,
transparency, and image-processing steps mimic experimental acquisition, it can
expose errors that are invisible in clean FOM-to-ROM tests: segmentation bias,
edge-location uncertainty, coordinate calibration errors, and systematic
distortion of the observed free surface. These effects can be included in the
Bayesian model either as larger measurement noise or as a correlated model-bias
term.


\FloatBarrier
\section*{Declaration of generative AI and AI-assisted technologies in the manuscript preparation process}

During the preparation of this work the author(s) used Anthropic's Claude in order to improve the language and grammar of the manuscript. After using this tool/service, the author(s) reviewed and edited the content as needed and take(s) full responsibility for the content of the published article.

\section*{Data and code availability}
The snapshot data, the reduced-order model, and the scripts that generate every figure are available at \url{https://github.com/parajal/bayesian-die-swell-main}.
\section*{Acknowledgments}
The authors thank M.A. Hulsen at the Eindhoven University of Technology (TU/e) for access to the TFEM software libraries.

\FloatBarrier

\bibliographystyle{elsarticle-num}
\bibliography{references}

\end{document}
